\scanpage{116}
\mspage{107}
% archive sheet number: 116 (lower left)
% L1-2
Ainsi, le théorème est établi. On en conclut que \deleted{\unsure{tout}}

% L3-13: the following block has long diagonal cancellation strokes.
\deleted{\uline{Corollaire} (\unsure{détermination|determination} des supports constructibles de $\unsure{M|\mathcal{M}}$). \deleted{Soit $X$}}

\deleted{1) Tout support constructible peut s'obtenir de la façon suivante : on prend une figure \deleted{\unsure{$f$}} $F=F_D$, $D=(T_1<_{\mathrm{pos}}T_2<\cdots<_{\mathrm{pos}}T_n)\in\operatorname{Drap}^{*}(\operatorname{Drap}^{*}(\unsure{\mathcal{L}|L}))$, et pour tout $T_i$, on prend une partie $\beta_i\subset\unsure{\partial|\delta}F_{T_i}=\partial T_i$. \deleted{\unsure{On lui}} On lui associe la partie constructible}
\[
\deleted{\Phi=\deleted{\unsure{F}}(F_D\setminus F_D\cap M_0)\cup\bigcup\beta_i=\bigcup_i((F_{T_i}\setminus F_{T_i}\cap M_0)\cup\beta_i)}
\]
\deleted{et}
\[
\deleted{\operatorname{supp}^{\circ}\Phi=\operatorname{Omb}^{\circ}(\Phi)}
\]
% L14-20: below the cancellation strokes.
Alors
\[
\operatorname{supp}^{\circ}\Phi=\operatorname{supp}^{\circ}\Phi'
\]
$\Phi'$ \unsure{déduit} de $\Phi$ en ``\unsure{recollant}'' les morceaux'', et ayant la propriété d'être \uline{``\unsure{non morcelée|non morcelé}''}
i.e.
\[
\forall x\in\Phi'\cap M_0,\quad\operatorname{card}\{X\in\Phi'\mid x\mathrel{\triangleleft}X\}\ne2
\]
(\uline{NB} Ce cardinal est \unsure{tjrs} $\in\{0,1,2\}$)

Ceci posé, on trouve.
